\section{Conclusions}
\label{sec:conclusions}

This study set out to close a practical blind spot: reproducible morphometrics can
say what kind of fabric a place is, and resilience assessment can say where risk is,
but the two rarely meet at the street and neighbourhood scale where adaptation is
designed. The proposed workflow binds them in one commensurable unit: a
250\,m grid cell that carries tessellation-cell morphometrics, network reach and
hazard exposure together. It converts them through an explain-then-optimize
sequence into transparent, fabric-specific adaptation priorities. The
explainable model diagnoses why fabrics are hot and which
intervention each needs, while the optimizer ranks where to act, with the
spatial-statistics and equity layer made explicit.

The four research questions close as follows. \textbf{RQ1:} yes, the seven
provisional strata are sharply differentiated on measured form, access and heat
(Table~\ref{tab:fabric-profiles}), although they consolidate into four
morphometric super-types. \textbf{RQ2:} stability is partial: the broad gradient
persists at 500\,m, but ARI = 0.38 and the loss of the finest mode show that
boundary membership remains scale-sensitive (Figure~\ref{fig:dendrogram}).
\textbf{RQ3:} the dominant mechanism is fabric-specific; coastal distance leads
globally, while large building footprints are the leading morphological warming
term in industrial/logistics and green cover leads several residential strata
(Figure~\ref{fig:shap}). \textbf{RQ4:} no fabric-level stratum is dominated, but
the cell-level screen identifies a selective frontier of 223 cells, and the
waterfront remains the highest mean-rank priority once exposure and vulnerability
are considered jointly (Table~\ref{tab:priority}).

A full-census application over the \.{I}zmir functional urban region, built entirely on open and official
open-data layers, demonstrates the chain end-to-end. Fabric is sharply
differentiated on measured form and consolidates into four morphometric
super-types whose broad structure persists across grid resolution (ARI 0.38, fair
agreement), with the finest mode resolution-sensitive.
Measured summer heat is highest in the coarse-grain industrial fabric, and although
the explainable model shows that, at metropolitan scale, form is a secondary
control on land-surface temperature behind the coastal gradient, morphology carries
modest but genuine skill ($R^2\approx0.09$ alone, $\approx0.27$ with context) and
building-footprint grain is the leading morphological mechanism. Most consequentially, coupling heat,
access, coastal exposure and age-based social vulnerability inverts a heat-only
reading: the coolest fabric, the waterfront, is the highest adaptation priority
because it is the most demographically vulnerable and the most exposed. This is exactly the
case that coarse, single-hazard assessment misses.
Planned extensions include seasonal and diurnal LST decomposition, pixel-scale SHAP
attribution, and multiscale geographically weighted regression for local
non-stationarity. These will deepen the mechanistic reading, and deployment on a second
coastal city will test inter-city transferability of the workflow and the
coastal-context finding.

The results are deliberately reported as screening, with their uncertainty and
weight-robustness, and the fabric strata remain provisional pending planning and
urban-design refinement. With frozen software, logged parameters and released code,
the workflow is transferable to other coastal metropolitan regions. It carries the
lesson that, in maritime cities, contextual drivers must travel alongside the
morphological ones.
